\documentclass[11pt]{article}
\usepackage{problem_set_style}
\usepackage{kotex}
\usepackage{booktabs}
\usepackage{array}
\begin{document}
\heading{}{샘플}{}{Spring 2027}{프로젝트 2 · 동료평가 샘플 (17주차 제출)}

\noindent{\small 아래는 \textbf{가상의 팀}이 쓴 샘플입니다. 수치는 형식을 보이기 위한 예시이며 실제 실험 결과가 아닙니다. 형식과 분량의 기준으로만 참고하세요.}

\noindent\textbf{평가 대상}: 다른 팀의 ``영화 리뷰 감성 분류'' 보고서

\begin{center}\small
\begin{tabular}{@{}p{3.4cm}p{9cm}@{}}\toprule
항목 & 평가 \\\midrule
주장의 명확성 & LSTM이 TF-IDF보다 낫다는 주장이 분명하다. \\
비교 대상 & TF-IDF + 로지스틱회귀 베이스라인이 없다. 왜 딥러닝인가가 정확도가 높아서 한 줄이다. \\
확장성 & 영화 리뷰만 다뤘다. 다른 도메인(상품 리뷰)에서의 확인이 없다. \\
재현성 & 데이터 출처, 시드, 에포크 수, 실행 명령이 모두 있어 다시 해 볼 수 있다. 좋다. \\
한계 & 긴 문장에서 성능이 떨어지는 것을 보이고 원인을 진단했다. 좋다. \\
\bottomrule\end{tabular}\end{center}

\noindent\textbf{반박 가능한 질문.} TF-IDF + 로지스틱회귀 같은 단순한 베이스라인을 돌렸나요? 그 정확도와 LSTM의 차이가 계산 비용을 감당할 만큼인가요?

\end{document}
